\section{Verification, integration, and further research}
\label{future:sec}

\subsection{What has been proved and what has been checked}

The proofs in this article establish the identities. Their main
analytic steps are normal convergence after pairing Fourier
coefficients with a smooth test function, explicit Taylor subtraction
at a singular endpoint, joint meromorphic continuation before taking
parameter coefficients, and finite Gamma/Bernoulli coefficient
algebra. None of the theorems is inferred from a numerical fit.

The most delicate points were checked independently of their initial
derivations. These include the constant Fourier mode, the signs of
the covariant Dirac derivatives, the minus sign on the wrapped
half-twist interval, the harmonic Mellin remainder at a Gamma zero,
the depth-mixing reflection shift \(a+u\), and the terminal
\(\mathcal A_{-1}=1/\Gamma(1+u)\) term in rational multiplication.
The full mathematical review is supplied in
\path{verification/independent_review.md}. This is an independent
derivation within the present research workflow, not a claim of
external peer review or proof-assistant formalization.

The exact and numerical replays have different roles.
Exact symbolic checks verify finite polynomial consequences of the
stated generating identities. Numerical checks compare independent
analytic representations, with endpoint cancellations removed before
evaluation. They are floating-point diagnostics, not rigorous
interval enclosures.

\begin{center}
\begin{tabularx}{\textwidth}{@{}p{0.29\textwidth}X@{}}
\toprule
Verification component & Recorded coverage\\
\midrule
Untwisted coefficient algebra &
Six explicit Stieltjes products; reflected phase coefficients through
degree eight for both signs; 128 polygamma product cases and
30 repeated-primitive recurrences.\\[3pt]
Harmonic exact algebra &
167 checks: local kernels, Hurwitz--Bernoulli reductions, residue
differentiation, reflection, primitives, multiplication for moduli
two and three, and all displayed specializations.\\[3pt]
Periodic quadrature &
15 checks at 50 decimal digits, including finite-part test-function
pairings and ordinary digamma/log-Gamma convolutions.\\[3pt]
Harmonic Mellin continuation &
Three regular-value checks and four Laurent cases at each of
\(\varepsilon=10^{-4},10^{-5},10^{-6}\), using an independent scalar
continued Mellin evaluator at 110 decimal digits.\\[3pt]
Half-twist and source repair &
Nine checks at 60 decimal digits: the convergent half-twist integral
at three shifts, its quarter-Gamma evaluation, four branch-corrected
arctanh integrals, and the exact spurious imaginary term.\\[3pt]
Twisted spectral normalization &
The separate twisted replay tests rational Lerch conversion,
the retained zero mode, and the covariant contact coefficients.
Its complete parameters and residuals are recorded with the script.\\
\bottomrule
\end{tabularx}
\end{center}

For the harmonic continuation, the scalar integrand is
\[
 \frac{e^{-(a+1)t}}{1-e^{-t}}\,
       \frac{[-\log(1-e^{-t})]^r}{r!}.
\]
The independent evaluator integrates its convergent local
logarithmic series on \((0,1/2)\), after analytic continuation of the
individual Mellin monomials, and performs direct quadrature on
\((1/2,\infty)\). It does not evaluate through the
\(\mathcal A_m\) coefficient theorem. Degrees 80 and 100 agreed to
better than \(1.7\times10^{-92}\) in the recorded tests. After
subtracting the asserted principal and constant terms, the remainder
was \(O(\varepsilon)\) in every tested case. That agreement diagnoses
the coefficient extraction; it does not certify a truncation bound
for arbitrary inputs.

The extra half-twist quadrature uses only digamma and polygamma
functions on the integral side, and Hurwitz Stieltjes constants or
\(\log\Gamma(1/4)\) on the comparison side. A Taylor quotient removes
the cancellation at \(x=0\). Its largest recorded scaled residual,
including the branch-audit checks, is below \(4\times10^{-61}\).
Here a scaled residual means
\[
 \frac{|\hbox{computed left side}-\hbox{computed right side}|}
      {1+|\hbox{computed right side}|}.
\]
A tiny residual is useful implementation evidence, not an equality proof.

\subsection{Suggested repository integration}

The package is designed for archival intake first and mathematical
integration second. The supplied source correction patch applies
against the pinned chapter bytes. The main article has modular
sources and an equivalent standalone source, plus complete replay
scripts and recorded results.

The integration should preserve the following distinctions:
\begin{enumerate}[leftmargin=1.6em]
\item Cite the existing untwisted continuations instead of replacing
      their provenance with the present rederivation.
\item Add the nonintegral-twist section as a response to the archived
      half-twist question, retaining its cutoff definition, covariant
      derivative, unit \(\delta_0\), and reflection modulation.
\item Add the shifted harmonic Laurent theorem together with its
      precise outer-shift convention. Moving the inner harmonic
      denominators as well defines a different family.
\item Keep each finite-part identity attached to its domain and
      normalization. An isolated ordinary formula does not determine
      its endpoint distribution.
\item Apply the two narrowly scoped source corrections separately
      from accepting the new research statements.
\end{enumerate}
The machine-readable claim ledger records proved additions,
consolidated antecedents, corrections, and unresolved targets.
The proposed patch does not change the status of \(S_6\) or \(S_8\).

\subsection{Further research questions}

The following questions concern exact identities and reductions.
They are proposals for further work, not claims supported merely by
high-precision numerical evidence.

\begin{question}[Products with genuinely different twists]
\label{future:q:unequaltwists}
For real nonintegral \(\theta,\eta\), with
\(\theta-\eta\notin\mathbb Z\), determine a useful finite or
convergently evaluable normal form for
\[
 \widehat{\mathcal K_{\alpha,\beta}^{\theta,\eta}}(n)
 =(2\pi i(n+\theta))^\alpha
  (2\pi i(n+\eta))^\beta .
\]
Which parameter configurations permit reduction to finitely many
single-twist Lerch jets, and when is a genuinely two-shift
special function necessary?
\end{question}

The equal-twist case is now solved by the Gamma quotient.
The half-twist involution reverses frequencies within one twist;
it does not solve the unequal-twist problem. A successful extension
should retain its Laurent endpoint terms and rational-character
coordinates, rather than only its nonsingular Fourier expression.

\begin{question}[Regular Laurent coefficients beyond the finite part]
\label{future:q:regularcoefficients}
For \(j\geq1\), determine exact representations and reduction rules
for
\[
 [\varepsilon^j]E_r(-m+\varepsilon;a).
\]
At rational \(a\), when can these coefficients be expressed using
ordinary Hurwitz or Dirichlet-\(L\) derivatives, and when do new
height-one or colored multiple-zeta derivatives remain?
\end{question}

The Gamma zero eliminates the analytic Mellin remainder through
the constant coefficient, but not beyond it. Thus the polynomial
theorem proved here gives the full singular and finite data without
settling the next regular coefficient. The independent scalar
Mellin representation supplies a rigorous starting point for
such a study.

\begin{question}[Lifting the reflection and multiplication laws]
\label{future:q:lift}
Find exact formulas for the analytic defects obtained when
\eqref{harm:eq:FPreflection} and \eqref{harm:eq:multiplication}
are lifted from Laurent coefficient polynomials to the full
meromorphic family \(E_r(s;a)\).
Can the defects be organized as convergent iterated integrals or
colored polylogarithms with a natural depth filtration?
\end{question}

The finite polynomial identities are exact, not asymptotic.
Nevertheless their proof uses only the local Mellin kernel. A
global lift must also control the rest of the integral. Keeping
that local/global separation explicit should prevent a false
unmodified Hurwitz multiplication rule for the harmonic family.

\begin{question}[Rational-twist arithmetic at singular orders]
\label{future:q:characters}
Combine the all-index twisted table with the manuscript's
conductor-sensitive character conversion. Determine economical
primitive-character formulas for the resulting rational-shift
convolutions, including every conductor logarithm and endpoint
term. Which combinations descend to a proper sub-conductor?
\end{question}

A finite coordinate algorithm is already available; the remaining
problem is its exact arithmetic simplification. Reduction to a
chosen spanning set is a meaningful target even when independence
of the evaluated constants is unknown.

\begin{question}[Translated cubic and higher log-Gamma products]
\label{future:q:moments}
With \(g(x)=\log\Gamma(x)-\tfrac12\log(2\pi)\), obtain exact reductions
for the two-parameter ordinary product
\[
 C_3(a,b)=\int_0^1g(x)g(\{x+a\})g(\{x+b\})\,dx,
 \qquad 0<a<b<1,
\]
and for higher moments beyond the known unshifted cubic formula.
At rational \(a,b\), which colored multiple-zeta derivatives remain
after exact character and symmetry reductions?
\end{question}

The unshifted cubic moment is already evaluated in the pinned chapters
\path{07-loggamma-moments.tex} and \path{07-tornheim-evaluation.tex},
and must not be proposed as an unsolved problem
\cite{RepoManuscript}. The target here retains two distinct translation
parameters. These are pointwise products, distinct from convolution
powers: Fourier multiplication introduces several independent indices.
A concrete starting family is the meromorphic continuation of
\[
 \mathcal T(s,t,u;z,w)=
 \sum_{m,n\geq1}\frac{z^mw^n}{m^s n^t(m+n)^u},
\]
and its order derivatives at the integer triples arising from
Kummer's Fourier coefficients, with phases fixed by \(a,b\).
The research target is an explicit
reduction theorem, not a numerical assertion that a selected
constant basis is complete or independent.

\begin{question}[Collision regularizations across different products]
\label{future:q:collisions}
When translated singularities collide, compare the finite part
obtained from the meromorphic convolution family with finite parts
obtained by coincident-point cutoff procedures. Determine the
exact local counterterm relating these prescriptions, and which
associativity or differentiation identities each preserves.
\end{question}

The distribution convolution itself is defined on the compact
circle. Evaluating it at its singular point is a different
operation and needs a separate prescription. The zero-mode
degeneration theorem controls a twist resonance, not every
possible collision regularization of pointwise products.

\begin{question}[The remaining Gaussian candidates]
\label{future:q:targets}
Can the level-four iterated-integral and conductor-jet reductions
in the manuscript be completed into an exact certificate for
\texttt{cycloquot:conj:S6} or \texttt{s8new:conj:S8}?
If a conjectured coefficient vector is false, can it be rejected
with a certified error enclosure and replaced by a precisely
stated new target?
\end{question}

The source already distinguishes these candidates from the proved
\(S_2,S_4\) identities and the rejected earlier \(S_8\) vector.
The present twist and harmonic identities should be used as
additional transformations, without suggesting that their
existence proves either outstanding relation.

\paragraph{A practical formalization programme.}
A useful accompanying project is to formalize the finite coefficient
algebra separately from the analytic continuation lemmas. The exact
Bell, Bernoulli, reflection, multiplication, and residue recurrences
give compact symbolic certificates at any requested depth. The
analytic layer should then discharge the local-uniform convergence,
endpoint subtraction, and uniqueness hypotheses once, instead of
replacing them by growing collections of numerical examples.

